\appendix

\section{Definitions and Notation}

A \emph{category} $\cC$ consists of objects, arrows $u:c\to d$ with composition $vu$ and identities $\id_c$. A \emph{functor} $F:\cC\to\cD$ preserves identities and composition. A \emph{natural transformation} $\eta:F\Rightarrow G$ between functors $F,G:\cC\to\cD$ is a family $\eta_c:Fc\to Gc$ with $\eta_dF(u)=G(u)\eta_c$ for $u:c\to d$. A \emph{presheaf} on $\cC$ is a functor $\cC^{\op}\to\Set$, and the presheaf category $[\cC^{\op},\Set]$ has natural transformations as morphisms. The constant presheaf on a set $X$ is written $\Delta X$. Epimorphisms and monomorphisms of presheaves are the componentwise surjective and injective natural transformations, and every natural transformation factors as an epimorphism followed by a monomorphism, the factorization being unique up to unique isomorphism.

A functor is \emph{full} if surjective on hom-sets, \emph{faithful} if injective on hom-sets, and \emph{essentially surjective} if every object is isomorphic to one in the image. An \emph{equivalence of categories} is a functor with a quasi-inverse, and Theorem~\ref{thm:equiv} characterizes it. A preorder is \emph{skeletal} if antisymmetric. The \emph{Cauchy completion} of a small category is obtained by splitting idempotents.

For a presheaf $P$ on $\cC$, a \emph{sieve} on $c$ is a set of arrows into $c$ closed under precomposition, a \emph{Grothendieck topology} selects covering sieves subject to the axioms of maximality, stability under pullback, and transitivity, and the sheaf condition for a covering sieve is that every matching family has a unique amalgamation. The \emph{limit} $\lim P$ of a presheaf viewed as a diagram is the set of compatible families $(s_c)_c$ with $P(u)(s_d)=s_c$.

An \emph{equivalence relation} on $X$ is a reflexive, symmetric, transitive relation, and the set $\mathrm{Eq}(X)$ of such relations is a complete lattice under inclusion. The \emph{kernel} of a map $g$ is $\{(x,y):g(x)=g(y)\}$. A \emph{labeled transition system} is $(S,\Sigma,\to)$ with $\to\subseteq S\times\Sigma\times S$, a \emph{simulation} from $S$ to $S'$ is a map or relation sending each transition to a transition, and a \emph{bisimulation} is a relation $R$ such that both $R$ and $R^{-1}$ are simulations. Partial maps are written $\pfn$, and a partial map $g$ is \emph{defined at} $x$ if $x\in\dom g$.

\section{Finite Examples and Counterexamples for Visibility, Disclosure, and Access Control}

\subsection*{B.1 Commutation of navigation and disclosure}

Let $t=(c,\texttt{go},c')$ and $S=\{0,1\}$ (or $\{0,1\}^2$ below). Entries give the configuration reached, and a dash-free ``undefined'' marks the failure of a composite to be defined.

\begin{center}
\begin{tabular}{llll}
\toprule
Setting & $(c,s)$ & $N_tD_f(c,s)$ & $D_fN_t(c,s)$\\
\midrule
global toggle, $\rho_t=\id$ & $(c,0)$ & $(c',1)$ & $(c',1)$\\
per-folder toggle, $S=\{0,1\}^2$, $\rho_t=\id$ & $(c,(0,0))$ & $(c',(1,0))$ & $(c',(0,1))$\\
reset on entry, $\rho_t\equiv0$, toggle & $(c,s)$ & $(c',0)$ & $(c',1)$\\
virtual target, $f_{c'}$ undefined & $(c,0)$ & $(c',1)$ & undefined\\
\bottomrule
\end{tabular}
\end{center}

The four rows are the cases of Proposition~\ref{prop:commute}: the first satisfies $\rho_tf_c=f_{c'}\rho_t$ and the other three violate it.

\subsection*{B.2 Authority and the three orders}

\begin{center}
\begin{tabular}{lllll}
\toprule
Context & Working directory & Spawned by & $\cA$ & Delegation\\
\midrule
$u$ & \texttt{/home} & (login) & $\{\mathrm{read},\mathrm{write}\}$ & top\\
$l$ & \texttt{/home/sandbox} & $u$ & $\{\mathrm{read}\}$ & $l\sqsubseteq u$\\
$s$ & \texttt{/home/bin} & $u$ & $\{\mathrm{read},\mathrm{write},\mathrm{admin}\}$ & none\\
\bottomrule
\end{tabular}
\end{center}

Containment gives $l\preceq u$ and $s\preceq u$, spawn gives $l$ and $s$ as children of $u$, and $\cA$ is monotone for delegation alone.

\subsection*{B.3 State-dependent accessibility}

The system $S=\{A,B\}\times\{0,1\}$ with initial state $(A,0)$, transitions $(A,0)\xrightarrow{\texttt{take}}(A,1)$ and $(A,1)\xrightarrow{\texttt{go}}(B,1)$, and $q(r,k)=r$ has context graph with edges $A\xrightarrow{\texttt{take}}A$ and $A\xrightarrow{\texttt{go}}B$. From the initial context $A$, the quotient admits the run \texttt{go}, which the system does not from $(A,0)$.

\subsection*{B.4 Shadowing and coherence}

Take states $(x_p,x_q)\in\{0,1\}^2$ for an outer scope $p$ and an inner scope $q$ that rebinds the same name. The flat observations $o_p=x_p$ and $o_q=x_q$ have kernels with classes $\{(0,0),(0,1)\},\{(1,0),(1,1)\}$ at $p$ and $\{(0,0),(1,0)\},\{(0,1),(1,1)\}$ at $q$, which are incomparable, so coherence fails for the arrow $p\to q$. The frame observations $o_p=x_p$ and $o_q=(x_p,x_q)$ have kernels with two classes at $p$ and four at $q$, nested as required.

\section{Worked Translations Among Filesystem, Lexical-Scope, and Transition-System Models}

Let $T$ be the rooted tree with root $R$, children $A$ and $B$, and a node $C$ whose parent is $A$, so that $C\preceq A\preceq R$ and $B\preceq R$. Let every local set be $L(p)=\{0,1\}$, so that $X=\{0,1\}^4$ has $16$ states with coordinates $(x_R,x_A,x_B,x_C)$.

\paragraph{Filesystem as region system.} The context category is $\cC^{\downarrow}$ with arrows $C\to A\to R$ and $B\to R$. The observation spaces are the products over nested nodes, and the kernels are as follows.
\begin{center}
\begin{tabular}{lccc}
\toprule
Context $p$ & $\Omega^{\downarrow}(p)$ & Classes of $\sk^{\downarrow}_p$ & Agree on\\
\midrule
$R$ & $\{0,1\}^4$ & $16$ & $R,A,B,C$\\
$A$ & $\{0,1\}^2$ & $4$ & $A,C$\\
$B$ & $\{0,1\}$ & $2$ & $B$\\
$C$ & $\{0,1\}$ & $2$ & $C$\\
\bottomrule
\end{tabular}
\end{center}
The family $\{R\}$ is generating, and $X_{\{R\}}=X$ has $16$ classes.

\paragraph{Lexical scope as accumulation system.} The context category is $\cC^{\uparrow}$ with arrows $R\to A\to C$ and $R\to B$, and the observation spaces are the products over ancestors.
\begin{center}
\begin{tabular}{lccc}
\toprule
Context $q$ & $\Omega^{\uparrow}(q)$ & Classes of $\sk^{\uparrow}_q$ & Agree on\\
\midrule
$R$ & $\{0,1\}$ & $2$ & $R$\\
$A$ & $\{0,1\}^2$ & $4$ & $R,A$\\
$B$ & $\{0,1\}^2$ & $4$ & $R,B$\\
$C$ & $\{0,1\}^3$ & $8$ & $R,A,C$\\
\bottomrule
\end{tabular}
\end{center}
The family $\{C\}$ alone is not generating, because $B$ is not a restriction of $C$, and $X_{\{C\}}$ has $8$ classes, identifying states that differ only in $x_B$. The family $\{B,C\}$ of innermost scopes is generating, and $X_{\{B,C\}}=X$ has $16$ classes, as Theorem~\ref{thm:coarsest}(c) predicts. The two kernel families are not related by any isomorphism of the indexing categories, since $\cC^{\downarrow}$ has one maximal element and two minimal ones, and $\cC^{\uparrow}$ has two maximal and one minimal.

\paragraph{Game world as transition system.} Let the configurations be $S=\{R,A,B,C\}\times\{0,1\}^4$, with actions $\texttt{down}_A$ from $R$ to $A$, $\texttt{down}_B$ from $R$ to $B$, $\texttt{down}_C$ from $A$ to $C$, \texttt{up} returning to the parent, and \texttt{flip} negating the bit of the current room, and let $q(r,x)=r$. The context graph $\Gamma_q$ has the edges $R\rightleftarrows A$, $R\rightleftarrows B$, and $A\rightleftarrows C$, which is the tree of $T$ with a return edge for every edge. The observation of a player in room $r$ is $x_r$, with discrete refinement category, so the tree appears only in the navigation structure.

\paragraph{What each model retains.}
\begin{center}
\begin{tabular}{lccc}
\toprule
 & Filesystem & Lexical scope & Game world\\
\midrule
role of the tree & refinement arrows $\preceq$ & refinement arrows $\succeq$ & navigation edges\\
observation at $p$ & subtree & ancestor chain & current room\\
generating contexts & root & leaves & none\\
state-dependent access & mode bits & none & guards\\
$\delta$ & \texttt{cd}, \texttt{chmod}, \texttt{rm} & binding and call & movement, \texttt{flip}\\
\bottomrule
\end{tabular}
\end{center}
The decorated tree $(T,|L|)$ with $|L|\equiv2$ is shared by the first two and supplies the state space of the third. The translation from the filesystem to the lexical-scope system that is the identity on states is a morphism of observational reducts only after reversing the context category, and it does not preserve $\delta$.

\section{Formal Inventory: Propositions, Counterexamples, and Open Questions}

The inventory lists every numbered result with its level: \emph{established} for results taken from the literature, \emph{proved here} for results whose statement or proof is supplied in this essay, and \emph{conjecture} or \emph{question} for open items. The final column gives the corresponding item of the proof spine, T for theorem, P for proposition, D for definition, where the spine item was incorporated.

\begin{center}
\small
\begin{longtable}{lllp{5.3cm}l}
\toprule
No. & Result & Level & Content & Spine\\
\midrule
\endfirsthead
\toprule
No. & Result & Level & Content & Spine\\
\midrule
\endhead
\ref{prop:treeposet} & Prop. & established & rooted trees as posets & \\
\ref{prop:catfromctx} & Prop. & established & categories from preorders and graphs & \\
\ref{prop:nomonotone} & Prop. & proved here & no forced monotonicity of $K,V,A$ & \\
\ref{def:crs} & Def. & definition & contextual representation system & D1\\
\ref{thm:refinement} & Thm. & proved here & observational refinement & T1\\
\ref{prop:authindep} & Prop. & proved here & independence of authority & P2\\
\ref{prop:commute} & Prop. & proved here & commutation criterion & P4\\
\ref{def:disclosure} & Def. & definition & disclosure refinement & D2\\
\ref{thm:disclosure} & Thm. & proved here & disclosure refinement & T3\\
\ref{prop:cdsubchmod} & Prop. & proved here & distinctness of \texttt{cd}, subshell, \texttt{chmod} & \\
\ref{def:scopegraph} & Def. & definition & scope graph, resolution policy & D3\\
\ref{thm:determinacy} & Thm. & proved here & determinacy of resolution & T5\\
\ref{prop:shadow} & Prop. & proved here & shadowing non-injective & P6\\
\ref{prop:visibility} & Prop. & proved here & visibility under nesting & \\
\ref{prop:graphgen} & Prop. & proved here & scope graphs generalize trees & P7\\
\ref{prop:authorders} & Prop. & proved here & authority and the three orders & \\
\ref{prop:overapprox} & Prop. & established & context graphs over-approximate & P18\\
\ref{prop:modal} & Prop. & proved here & modal return not a graph edge & P18\\
\ref{prop:flat} & Prop. & proved here & flat environments violate coherence & \\
\ref{prop:gluing} & Prop. & established & gluing criterion & T10\\
\ref{prop:disagree} & Prop. & proved here & disagreement excludes realization & \\
\ref{ex:triangle} & Ex. & proved here & obstruction to gluing & T10 target\\
\ref{ex:parity} & Ex. & proved here & ambiguity, failure of separation & T10 target\\
\ref{thm:canonical} & Thm. & established, applied & canonical presheaf, universal property & T8\\
\ref{prop:globalindist} & Prop. & proved here & global indistinguishability & P9\\
\ref{thm:coarsest} & Thm. & proved here & coarsest quotient for a family of contexts & \\
\ref{prop:Qglue} & Prop. & proved here & gluing criterion for $Q$ & \\
\ref{prop:minimization} & Prop. & established & behavioral minimization & \\
\ref{def:translation} & Def. & definition & translation & D4\\
\ref{prop:crs} & Prop. & proved here & composition of translations & P0\\
\ref{thm:presindist} & Thm. & proved here & translation preserves indistinguishability & T11\\
\ref{prop:reflection} & Prop. & proved here & reflection requires injectivity & P12\\
\ref{thm:central} & Thm. & proved here & common abstraction without equivalence & T13\\
\ref{prop:trivialabs} & Prop. & proved here & terminal system, trivial abstraction & \\
\ref{thm:terminalobs} & Thm. & established, applied & terminal observed system & \\
\ref{cor:retains} & Cor. & proved here & what the canonical abstraction forgets & \\
\ref{ex:flip} & Ex. & proved here & same abstraction, no translation back & P16\\
\ref{prop:orientation} & Prop. & proved here & orientation duality & P17\\
\ref{thm:equiv} & Thm. & established & equivalence criterion & T15\\
\ref{prop:surjnotequiv} & Prop. & proved here & surjectivity is not equivalence & P14\\
\ref{prop:nocollapse} & Prop. & proved here & no unrestricted faithful collapse & P20\\
\ref{prop:posetrigid} & Prop. & proved here & equivalence of posets is isomorphism & \\
\ref{cor:morita} & Cor. & established & Morita rigidity for posets & \\
\ref{thm:relequiv} & Thm. & proved here & relative equivalence & \\
\ref{prop:commonquot} & Prop. & proved here & common quotient, no mutual translation & P16\\
\ref{prop:PAHA} & Prop. & established & classical and intuitionistic arithmetic & \\
\ref{prop:glivenko} & Prop. & established & Glivenko map & \\
\ref{q:interp} & Question & open & interpretations and preservation & \\
\ref{prop:fullabs} & Prop. & established & full abstraction as canonical quotient & \\
\ref{prop:reify} & Prop. & established & contexts migrate into data & \\
\ref{prop:univinsufficient} & Prop. & established & universality is insufficient & P19\\
\ref{thm:partrep} & Thm. & proved here & partition representation & \\
\ref{thm:treerep} & Thm. & proved here & representation of tree systems & \\
\ref{prop:functorial} & Prop. & proved here & functoriality of the canonical quotient & \\
\ref{q:pushout} & Question & open & pushouts of correspondences & \\
\ref{prop:reach} & Prop. & proved here & reachability under translation & \\
\ref{prop:authority} & Prop. & proved here & authority, prohibition, safety & \\
\ref{prop:reconstruction} & Prop. & proved here & reconstruction under translation & \\
\ref{q:complexity} & Question & open & complexity of translation & \\
\ref{prop:unfold} & Prop. & established & unfolding is a bisimulation & \\
\ref{prop:unfoldloss} & Prop. & proved here & information lost by unfolding & \\
\ref{prop:descent} & Prop. & proved here & descent & \\
\ref{ex:sharedroom} & Ex. & proved here & shared room, two routes & \\
\ref{conj:descent} & Conj. & conjecture & descent for dynamics & \\
\bottomrule
\end{longtable}
\end{center}

The three levels of claim announced in Section~1 are the three values of the third column. No result classified as established is claimed to be new, and the results classified as proved here are elementary in their proofs and are offered for their formulation and their use in organizing the comparison.

\section{Related Frameworks and Bibliographic Notes}

The categorical background is standard \cite{maclane,borceux}, with an accessible account in \cite{awodey} and the treatment of sheaves and sites in \cite{maclanemoerdijk}. The programming-language background draws on \cite{harper,pierce}, program analysis on \cite{nielson}, and the semantics of objects on \cite{abadicardelli}. The scope-graph formulation of name resolution adopted in Section~6 is that of \cite{scopegraphs}, and the use of presheaf-valued observation relative to contexts follows the spirit of sheaf-theoretic treatments of contextuality \cite{abramsky}. Action semantics \cite{mosses} and categorical combinators \cite{curien} are alternative frameworks for modular semantics of programming languages, and the groupoid interpretation of type theory \cite{hofmannstreicher} is a categorical semantics in which identity is itself structured. For translation among logics, the theory of institutions \cite{goguen} provides a general framework in which a logical system is a functor from signatures to satisfaction relations and translations are morphisms of such functors, and relating contextual representation systems to institutions is a natural next step noted in Section~20. For access control, the original formulations are \cite{lampson,dennisvanhorn}, the undecidability of safety is \cite{hru}, and the object-capability discipline is discussed in \cite{capmyths}. The suggested reading list from which the present development started is a list of starting points for a literature review, and the inclusion of a work in it is not a claim that the construction proposed here has been established there.

\paragraph{Comparison with institutions.} An institution assigns to each signature a category of models and a set of sentences, with a satisfaction relation that is invariant under change of signature \cite{goguen}. Three features of the present framework are representable in that setting, and three are not without extra structure. The observational kernel $\sk_c$ of a context is the theory-indistinguishability relation of the models relative to a fixed set of sentences, so the canonical presheaf $Q$ corresponds to the quotient of models by elementary equivalence relative to the sentences available at $c$, and the restriction maps correspond to signature inclusions. The terminal observed system is the set of global sections of the presheaf of sentence-evaluations. The translation conditions (O) are the satisfaction condition restricted to the observed sentences. What an institution does not carry is an order on contexts that is independent of signature inclusion, a component of authority separate from satisfaction, and a transition function on models; the lax condition (A) and the simulation condition (D) have no counterpart there, and the failure of the commutation criterion of Section~4 is invisible to an institution because navigation is not a morphism of signatures. The framework is therefore best described as institution-like in its observational reduct and as a strict extension of it in authority and dynamics, and a precise comparison, which would exhibit the observational reduct of $\CRS$ as a full subcategory of a category of institutions, has not been carried out here.

\paragraph{Comparison with scope graphs and coalgebraic semantics.} The scope-graph calculus of \cite{scopegraphs} defines name resolution by reachability under a path policy, which corresponds in this essay to the specification of the observation presheaf from a graph, and Proposition~\ref{prop:graphgen} records that scope graphs strictly generalize scope trees. The calculus does not address the comparison of resolution disciplines across languages, nor authority or transitions. Coalgebraic treatments of behavioral equivalence \cite{sangiorgi} characterize the coarsest bisimulation of a transition system, and Proposition~\ref{prop:minimization} is the instance of this for the partial transition function used here; the observational quotient $X_{\cU}$ differs from the behavioral quotient exactly by the distinction illustrated in Example~\ref{ex:flip}.

\paragraph{Simulation-based accounts of physical reasoning.} A research program in cognitive science proposes that intuitive physical judgment is produced by approximate probabilistic simulation, in the manner of a game physics engine, restricted to the objects and interactions relevant to the question posed \cite{ullman}. This is an empirical hypothesis about people and is not a consequence of anything proved here, and the mathematics of this essay offers no evidence for or against it. A formal observation is nevertheless available. A simulator that instantiates only the objects relevant to a query computes, for each query, an observation at a restricted context, and the requirement that answers to related queries cohere is the coherence condition $o_c=\Omega(u)\circ o_d$. Whether the sets of queries and simulations in such an account form a presheaf satisfying the separation and gluing conditions of Section~10, and whether the failure modes of that section occur in it, are questions that the experimental literature has not posed in this form, and they are recorded as Question~\ref{q:simulation}.

\begin{question}[Simulation-based accounts as contextual systems]\label{q:simulation}
Can a simulation-based account of physical judgment be presented as a contextual representation system whose contexts are the sets of objects instantiated for a query, and, if so, do the observation maps satisfy separation and gluing, or do the obstruction and ambiguity modes of Section~10 occur?
\end{question}

\begin{thebibliography}{99}
\bibitem{abadicardelli} M. Abadi and L. Cardelli, \emph{A Theory of Objects}, Springer, 1996.
\bibitem{abramsky} S. Abramsky and A. Brandenburger, The sheaf-theoretic structure of non-locality and contextuality, \emph{New Journal of Physics} 13 (2011), 113036.
\bibitem{aczel} P. Aczel, The type theoretic interpretation of constructive set theory, in \emph{Logic Colloquium '77}, North-Holland, 1978, 55--66.
\bibitem{awodey} S. Awodey, \emph{Category Theory}, 2nd ed., Oxford University Press, 2010.
\bibitem{barretthalvorson} T. W. Barrett and H. Halvorson, Morita equivalence, \emph{Review of Symbolic Logic} 9 (2016), 556--582.
\bibitem{borceux} F. Borceux, \emph{Handbook of Categorical Algebra 1: Basic Category Theory}, Cambridge University Press, 1994.
\bibitem{capmyths} M. S. Miller, K.-P. Yee, and J. Shapiro, Capability myths demolished, Technical Report SRL2003-02, Johns Hopkins University Systems Research Laboratory, 2003.
\bibitem{carneiro} M. Carneiro, The type theory of Lean, unpublished manuscript, 2019.
\bibitem{curien} P.-L. Curien, \emph{Categorical Combinators, Sequential Algorithms, and Functional Programming}, 2nd ed., Birkh\"auser, 1993.
\bibitem{dennisvanhorn} J. B. Dennis and E. C. Van Horn, Programming semantics for multiprogrammed computations, \emph{Communications of the ACM} 9 (1966), 143--155.
\bibitem{felleisen} M. Felleisen, On the expressive power of programming languages, \emph{Science of Computer Programming} 17 (1991), 35--75.
\bibitem{friedman} H. Friedman, Classically and intuitionistically provably recursive functions, in \emph{Higher Set Theory}, Lecture Notes in Mathematics 669, Springer, 1978, 21--27.
\bibitem{glivenko} V. Glivenko, Sur quelques points de la logique de M. Brouwer, \emph{Bulletin de la Classe des Sciences, Acad\'emie Royale de Belgique} 15 (1929), 183--188.
\bibitem{goguen} J. A. Goguen and R. M. Burstall, Institutions: abstract model theory for specification and programming, \emph{Journal of the ACM} 39 (1992), 95--146.
\bibitem{harper} R. Harper, \emph{Practical Foundations for Programming Languages}, 2nd ed., Cambridge University Press, 2016.
\bibitem{hofmannstreicher} M. Hofmann and T. Streicher, The groupoid interpretation of type theory, in \emph{Twenty-Five Years of Constructive Type Theory}, Oxford University Press, 1998, 83--111.
\bibitem{hru} M. A. Harrison, W. L. Ruzzo, and J. D. Ullman, Protection in operating systems, \emph{Communications of the ACM} 19 (1976), 461--471.
\bibitem{lambekscott} J. Lambek and P. J. Scott, \emph{Introduction to Higher Order Categorical Logic}, Cambridge University Press, 1986.
\bibitem{lampson} B. W. Lampson, Protection, \emph{Proceedings of the 5th Princeton Conference on Information Sciences and Systems}, 1971, 437--443; reprinted in \emph{ACM Operating Systems Review} 8 (1974), 18--24.
\bibitem{lawvere} F. W. Lawvere, Functorial semantics of algebraic theories, \emph{Proceedings of the National Academy of Sciences} 50 (1963), 869--872.
\bibitem{maclane} S. Mac Lane, \emph{Categories for the Working Mathematician}, 2nd ed., Springer, 1998.
\bibitem{maclanemoerdijk} S. Mac Lane and I. Moerdijk, \emph{Sheaves in Geometry and Logic}, Springer, 1992.
\bibitem{minamide} Y. Minamide, G. Morrisett, and R. Harper, Typed closure conversion, \emph{Proceedings of POPL 1996}, 271--283.
\bibitem{moore} E. F. Moore, Gedanken-experiments on sequential machines, in \emph{Automata Studies}, Princeton University Press, 1956, 129--153.
\bibitem{morris} J. H. Morris, Lambda-calculus models of programming languages, Ph.D. thesis, Massachusetts Institute of Technology, 1968.
\bibitem{mosses} P. D. Mosses, \emph{Action Semantics}, Cambridge University Press, 1992.
\bibitem{mro} K. Barrett, B. Cassels, P. Haahr, D. A. Moon, K. Playford, and P. T. Withington, A monotonic superclass linearization for Dylan, \emph{Proceedings of OOPSLA 1996}, 69--82; and M. Simionato, The Python 2.3 method resolution order, Python documentation.
\bibitem{nielson} F. Nielson, H. R. Nielson, and C. Hankin, \emph{Principles of Program Analysis}, Springer, 1999.
\bibitem{scopegraphs} P. Neron, A. Tolmach, E. Visser, and G. Wachsmuth, A theory of name resolution, in \emph{Programming Languages and Systems (ESOP 2015)}, Lecture Notes in Computer Science 9032, Springer, 2015, 205--231.
\bibitem{pierce} B. C. Pierce, \emph{Types and Programming Languages}, MIT Press, 2002.
\bibitem{plotkin} G. D. Plotkin, LCF considered as a programming language, \emph{Theoretical Computer Science} 5 (1977), 223--255.
\bibitem{posix} The Open Group, \emph{The Open Group Base Specifications, Issue 7} (IEEE Std 1003.1-2017), utility \texttt{cd}.
\bibitem{pyref} Python Software Foundation, \emph{The Python Language Reference}, section on the execution model: naming and binding.
\bibitem{ullman} T. D. Ullman, E. Spelke, P. Battaglia, and J. B. Tenenbaum, Mind games: game engines as an architecture for intuitive physics, \emph{Trends in Cognitive Sciences} 21 (2017), 649--665.
\bibitem{reynolds} J. C. Reynolds, Definitional interpreters for higher-order programming languages, \emph{Proceedings of the ACM Annual Conference}, 1972, 717--740.
\bibitem{rogers} H. Rogers Jr., G\"odel numberings of partial recursive functions, \emph{Journal of Symbolic Logic} 23 (1958), 331--341.
\bibitem{sangiorgi} D. Sangiorgi, \emph{Introduction to Bisimulation and Coinduction}, Cambridge University Press, 2011.
\bibitem{troelstravandalen} A. S. Troelstra and D. van Dalen, \emph{Constructivism in Mathematics: An Introduction}, vol.~1, North-Holland, 1988.
\bibitem{vpa} R. Alur and P. Madhusudan, Visibly pushdown languages, \emph{Proceedings of STOC 2004}, 202--211.
\end{thebibliography}
